% ===========================================================================
% PART 3 -- the five things an audit of the first draft found missing.
% Four are measured here; the fifth is a foundational assumption that had been
% used without being named.
% ===========================================================================
\section{Five things the first draft assumed}
\label{sec:gaps}

An audit of this paper against its own standard found five loads being carried by
statements that were never made. Four are now measured. The fifth is not
measurable and is instead named, because using it silently is the error.

\subsection{G0 --- interpersonal comparability, the assumption under
\texorpdfstring{§\ref{sec:agg}}{the aggregation section}}

Every aggregate in §\ref{sec:agg} adds one person's numbers to another's. That
operation is meaningless unless the numbers live on a common scale, and
\textbf{nothing in the elicitation supplies one}. Person $1$ writing $+8$ and
person $2$ writing $+8$ are two marks on two private rulers.

\step{What the paper actually did.} It min--max normalised each person to
$[0,1]$ (§\ref{sec:notation}), so that everyone's best option is worth $1$ and
worst worth $0$. That is a \emph{stipulation of equal stakes}: it asserts that
every participant cares equally much about the whole question.

\step{Why the stipulation is not innocent.} Consider two participants: one for
whom all four responses are nearly equivalent, and one for whom the choice is
grave. Min--max gives them identical influence over the aggregate. A different
and equally defensible stipulation --- normalise by a common scale factor, so
that intensity counts --- gives the second participant more. The two produce
different winners.

\begin{proposition}[The deficiency figure is scale-relative]
The Blackwell deficiency of $0.1298$ reported in §\ref{sec:agg} is a number about
the pair (data, interpersonal scale). It is not identified by the data alone,
because the data contain no information distinguishing ``this person's $+8$ means
what that person's $+8$ means'' from any other calibration.
\end{proposition}

\step{Why this is the deepest limitation in the paper.} Arrow's theorem is
precisely a theorem about what can be done \emph{without} interpersonal
comparison; Harsanyi's utilitarian representation requires it. So the choice made
here silently determines which side of a hundred-year-old divide the analysis
sits on. It is now stated. It cannot be resolved at this site: resolving it needs
an elicitation that measures intensity on a shared anchor, which is a design
requirement and appears in §\ref{sec:design}.

\mesoboundary{Every aggregate quantity in this paper --- the deficiency, the rule
comparisons, the maximin figures --- is conditional on min--max interpersonal
normalisation. The \emph{ordinal} results (Condorcet structure, IIA rate,
agreement rates) are not, because they never add across people.}

\subsection{G1 --- ties, asserted and now measured}

The paper's decision rules use $\argmax$, and exact ties would make them
ill-defined. The first draft asserted that ties occur with probability zero.
Measured over all $968$ prompts:

\begin{table}[H]\centering\small
\begin{tabular}{lr}
\toprule
exact ties between the top two scores & $\mathbf{0}$ of $968$\\
relative gaps below $10^{-6}$ & $0$ of $968$\\
\bottomrule
\end{tabular}
\end{table}

\step{Why it needed checking rather than asserting.} ``Probability zero'' is a
statement about a continuous model; the data are finite and produced by a
tie-prone construction (sums of a few dozen products of bounded numbers). The
claim happened to be true. It was not entitled to be assumed, and an assertion
that turns out true is still an assertion until it is run.

\subsection{G2 --- manipulability, and what it does to ``sincere''}

Every number in this paper reads the elicited weights as what people believe.
Gibbard--Satterthwaite guarantees that any non-dictatorial rule over three or
more alternatives \emph{can} be manipulated; it says nothing about how often, and
the rate is what decides whether the sincerity reading is safe.

\begin{protocol}[Crudest manipulation]
For each (annotator, prompt): compute the collective winner honestly, then
replace \emph{all} of that annotator's weights by $+10$, and separately by $-10$,
and recompute. Score both outcomes by that annotator's own stated ranking.
\end{protocol}

\begin{table}[H]\centering\small
\begin{tabular}{lr}
\toprule
(annotator, prompt) pairs tested & $15{,}482$\\
pairs where saturating the weights yields a \emph{strictly better} outcome & \\
\quad by that annotator's own ranking & $604 = \mathbf{3.9\%}$\\
\bottomrule
\end{tabular}
\end{table}

\begin{remark}[This is a lower bound, and that is the point]
The manipulation tested is the least intelligent one available --- push everything
to one end. A participant choosing which criteria to inflate can only do better.
So $3.9\%$ is a floor on the rate at which a participant had something to gain by
misreporting, and the sincerity reading of the weights is safe only to that
degree. No claim in this paper depends on sincerity being exact; several would
change in magnitude if the true manipulation rate were much larger, and that is
recorded in §\ref{sec:limits}.
\end{remark}

\subsection{G3 --- the judge scores criteria against the responses they were written about}

This is a circularity the first draft did not state. The criteria were written
\emph{after} the participant saw the four responses, typically to explain why one
of them is best. The judge then scores those criteria \emph{against those same
four responses}.

\begin{table}[H]\centering\small
\begin{tabular}{lr}
\toprule
criterion--author pairs & $95{,}553$\\
satisfaction of the author's own top-ranked response & $0.5831$\\
mean satisfaction of the other three & $0.5466$\\
difference & $+0.0365$ (se $0.0008$, $z=+46.2$)\\
\bottomrule
\end{tabular}
\end{table}

\step{How to read this, and how not to.} It is not a defect of the judge. It is
what ``I wrote down why $X$ is best'' means: a criterion authored by someone who
ranked $X$ first is, by construction, more satisfied by $X$. What it establishes
is that satisfaction and choice are \textbf{two views of one act}, not two
independent measurements of a latent value. Every correlation between the
criteria route and the ranking route --- including the $0.393$ of
\cref{tab:selfother} and the $68.4\%$ of \cref{tab:agree} --- is inflated by this
mechanism, and the inflation cannot be removed without response-blind criteria.

\subsection{G4 --- the two elicitations disagree beyond their own unreliability}

The criteria route and the ranking route select different winners on $31.6\%$ of
prompts (\cref{tab:agree}). A natural deflation is: both are noisy, and noise
alone explains it. \Cref{thm:atten} answers this exactly, because it says a
correlation measured through unreliable instruments is attenuated by
$\sqrt{\mathrm{rel}_1\mathrm{rel}_2}$ --- so dividing by that factor removes the
noise and leaves the disagreement.

\begin{table}[H]\centering\small
\begin{tabular}{lrr}
\toprule
& split-half & Spearman--Brown\\
\midrule
reliability of the ranking route & $0.7702$ & $0.8702$\\
reliability of the criteria route & $0.8882$ & $0.9408$\\
\midrule
observed cosine between the two routes & & $0.6265$\\
\textbf{disattenuated} $=0.6265/\sqrt{0.8702\times0.9408}$ & & $\mathbf{0.6924}$\\
\bottomrule
\end{tabular}
\label{tab:disatten}
\end{table}

\mesotakeaway{Correcting \emph{both} routes for their own unreliability raises
the agreement only from $0.626$ to $0.692$. The two elicitations still do not
coincide. So roughly $31\%$ of the gap between what people \emph{write} and what
they \emph{choose} is real disagreement, not measurement error --- and by G3 this
is an \emph{under}statement, because the circularity inflates the observed
agreement in the first place.}

\begin{corollary}[The paper's largest quantity, restated at the panel level]
\cref{tab:selfother} reported $0.393$ for one person against their own ranking.
\cref{tab:disatten} reports $0.692$ for the whole panel's criteria against the
whole panel's rankings, disattenuated. The two are consistent: averaging over
people removes individual noise (\cref{thm:sb}) and raises the number, but the
residual gap does not close. Whatever the pipeline does downstream, the largest single
discrepancy in the chain sits \emph{inside} the elicitation, between two things
the same people said in the same session.
\end{corollary}
